\documentclass[../main.tex]{subfiles}

\begin{document}

\ifSubfilesClassLoaded{

    %\tableofcontents
    
    %\setcounter{chapter}{}
}{}

\chapter{ Tangent Vectors }

\section{Tangent Space}

\subsection{Geometric Tangent Vectors}

\begin{definition}{}{}
     If $a$ is a point of $\mathbb{R}^n$, a map $w: C^\infty(\mathbb{R}^n) \to \mathbb{R}$ is called a \dfntxt{derivation at a} if it is linear over $\mathbb{R}$ and satisfies the following product rule:
$$w(fg) = f(a)wg + g(a)wf.$$
Let $T_a\mathbb{R}^n$ denote the set of all derivations of $C^\infty(\mathbb{R}^n)$ at $a$. Clearly, $T_a\mathbb{R}^n$ is a vector space under the operations
$$(w_1 + w_2)f = w_1f + w_2f, \quad (cw)f = c(wf).$$
\end{definition}

\begin{proposition}{}{}
Let $a \in \mathbb{R}^n$.
\begin{enumerate}
\item[(a)] For each geometric tangent vector $v_a \in \mathbb{R}_a^n$, the map $D_v|_a : C^\infty(\mathbb{R}^n) \to \mathbb{R}$ defined by\[
D_v|_a f = D_v f(a) = \frac{d}{dt}\Big|_{t=0} f(a+tv).
\]
  is a derivation at $a$.
\item[(b)] The map $v_a \mapsto D_v|_a$ is an isomorphism from $\mathbb{R}_a^n$ onto $T_a\mathbb{R}^n$.
\end{enumerate}
\end{proposition}

\subsection{Tangent Vectors on Manifolds}

\begin{definition}{}{}
    Let $M$ be a smooth manifold with or without boundary, and let $p$ be a point of $M$. A linear map $v: C^\infty(M) \to \mathbb{R}$ is called a \dfntxt{derivation at p} if it satisfies
\begin{equation}
    v(fg) = f(p)vg + g(p)vf \quad \text{for all } f,g \in C^\infty(M).
\end{equation}
The set of all derivations of $C^\infty(M)$ at $p$, denoted by $T_p M$, is a vector space called the \dfntxt{tangent space to M at p}. An element of $T_p M$ is called a \dfntxt{tangent vector at p}.
\end{definition}

\begin{lemma}{Properties of Tangent Vectors on Manifolds}{}
Suppose $M$ is a smooth manifold with or without boundary, $p \in M$, $v \in T_p M$, and $f, g \in C^\infty(M)$.
\begin{enumerate}
\item[(a)] If $f$ is a constant function, then $v f = 0$.
\item[(b)] If $f(p) = g(p) = 0$, then $v(fg) = 0$.
\end{enumerate}
\end{lemma}
\section{The Differential of a Smooth Map}

\subsection{Differential on Euclidean Space}

\begin{lemma}{}{}
    Let \(  U\subseteq \mathbb{R} ^{m}, V\subseteq \mathbb{R} ^{n}  \) are open sets. \(  f:U\to V  \) is smooth. Take \(  a \in v  \), define a map \[
  \begin{aligned}
    \left( df \right)_{a}:T_{a}U&\to T_{f\left( a \right) }V \\ 
     \left( df_{a} \right)\left( D_{v}^{a} \right)&\mapsto D_{w}^{f\left( a \right) },\quad w= J_{f}\left( a \right) v
  \end{aligned}
    \]  Then for each \(  g \in C^{\infty}\left( V \right)   \), we have \[
   \left(  \left( df_{a} \right)\left( v \right) \right) \left( g \right)   = D_{v}^{a}\left( g\circ f \right) 
    \] 
\end{lemma}
\begin{note}
    \(  \left( df_{a} \right)\left( v \right)    \)对应的导子恰好是在点 \(  f\left( a \right)   \)处将 \(  g\in C^{\infty}\left( V \right)   \)   映到 \(  D_{v}^{a}\left( g\circ f \right)   \)的导子. 
\end{note}

\begin{definition}{}{}
    If $M$ and $N$ are smooth manifolds with or without boundary and $F: M \to N$ is a smooth map, for each $p \in M$ we define a map
\[
dF_p: T_pM \to T_{F(p)}N,
\]
called the \dfntxt{differential of F at p} . Given $v \in T_pM$, we let $dF_p(v)$ be the derivation at $F(p)$ that acts on $f \in C^\infty(N)$ by the rule
\[
dF_p(v)(f) = v(f \circ F).
\]
\end{definition}

\begin{proposition}{Properties of Differentials}{}
Let $M$, $N$, and $P$ be smooth manifolds with or without boundary, let $F: M \to N$ and $G: N \to P$ be smooth maps, and let $p \in M$.
\begin{enumerate}
\item[(a)] $dF_p: T_p M \to T_{F(p)}N$ is linear.
\item[(b)] $d(G \circ F)_p = dG_{F(p)} \circ dF_p: T_p M \to T_{G \circ F(p)}P$.
\item[(c)] $d(\mathrm{Id}_M)_p = \mathrm{Id}_{T_p M}: T_p M \to T_p M$.
\item[(d)] If $F$ is a \dfntxt{diffeomorphism}, then $dF_p: T_p M \to T_{F(p)}N$ is an isomorphism, and $(dF_p)^{-1} = d(F^{-1})_{F(p)}$.
\end{enumerate}
\end{proposition}
\section{Tangent Bundle}
\begin{definition}{}{}
    Given a smooth manifold $M$ with or without boundary, we define the \dfntxt{tangent bundle of $M$}, denoted by $TM$, to be the disjoint union of the tangent spaces at all points of $M$:
$$TM = \bigsqcup_{p \in M} T_p M.$$
\end{definition}
\begin{remark}
\begin{itemize}
    \item \dfntxt{projection map} $\pi : TM \to M$, $\pi(p, v) = p$. 
\end{itemize}

\end{remark}

\begin{proposition}{}{}
    For any smooth \(  n  \)-manifold \(  M  \), then tangent bundle \(  TM  \) has a natural topology and smooth structure that make it into a \(  2n  \)-dimensional smooth manifold. With respect to this structure, the    projection \(   \pi: TM\to M  \) is smooth. 
\end{proposition}
\begin{proof}
    Given a chart \(  \left( U, \varphi  \right)   \) for \(  M  \), \(   \pi ^{-1} \left( U \right)\subseteq TM   \) is the set of all tangent vectors to \(  M  \) at all points of \(  U  \). Define a map \(   \tilde{\varphi} :  \pi ^{-1} \left( U \right)\to \mathbb{R} ^{2n}   \) by \[
     \tilde{\varphi} \left( v^{i}\left. \frac{\partial }{\partial x^{i}} \right|_{p} \right)= \left( x^{1}\left( p \right),\cdots ,x^{n}\left( p \right),v^{1},\cdots ,v^{n}   \right)  
    \]      
\end{proof}
\end{document}